% 03-research-wps.tex -- Research Areas (Work Packages)
\chapter{Research Areas}
\label{ch:research-wps}

The consortium's research is organized into five thematic work packages, each
addressing a distinct frontier of digital finance.  Every work package is led by
a beneficiary institution and draws on expertise from across the partnership.
Doctoral candidates are embedded in one primary work package while benefiting
from cross-WP collaboration, shared datasets, and joint publications.

% ====================================================================
% WP 1
% ====================================================================
\wpheader{1}{Towards a European Financial Data Space}{Babe\c{s}-Bolyai University (BBU)}

\subsection*{Research Motivation}

Researchers participating in WP1 will explore how Europe can build a shared,
privacy-preserving financial data ecosystem.  As financial institutions generate
ever-larger volumes of data, questions of access, quality, and governance become
critical.  This work package tackles the technical and regulatory challenges of
creating federated data infrastructures that comply with GDPR while enabling
innovation.

\subsection*{Key Research Questions}
\begin{itemize}[nosep]
  \item How can federated learning enable cross-border financial analytics
        without centralizing sensitive data?
  \item What synthetic data generation techniques best preserve statistical
        utility while guaranteeing privacy?
  \item How should data governance frameworks balance openness with regulatory
        compliance under GDPR and the Data Act?
  \item What role can privacy-enhancing technologies play in building trust among
        financial data providers?
\end{itemize}

\subsection*{Methodological Approach}

Researchers will combine federated learning architectures, differential privacy
mechanisms, and synthetic data generation (GANs, copula-based methods) with
empirical case studies drawn from consortium banking partners.

\subsection*{Expected Outcomes}
\begin{itemize}[nosep]
  \item Open-source federated learning toolkit for financial applications
  \item Benchmark datasets with validated synthetic counterparts
  \item Policy recommendations for a European Financial Data Space
  \item Publications in top-tier data science and finance venues
\end{itemize}

\subsection*{Related Training}
MST-01 (Foundation of Data Science), AST-01 (Synthetic Data Generation),
AST-02 (Anomaly Detection in Big Data), AST-03 (Natural Language Processing with Transformers),
AST-04 (Dependence Structures in High-Frequency Financial Data).

\subsection*{Contributing Partners}
\begin{center}
\partnerlogo{BBU}\hspace{3mm}%
\partnerlogo{WU}\hspace{3mm}%
\partnerlogo{ASE}\hspace{3mm}%
\partnerlogo{CAR}
\end{center}

\newpage

% ====================================================================
% WP 2
% ====================================================================
\wpheader{2}{AI for Financial Markets}{WU Vienna}

\subsection*{Research Motivation}

Researchers participating in WP2 will advance the application of modern
artificial intelligence to European financial markets.  While deep learning has
transformed many domains, its adoption in finance faces unique challenges:
non-stationarity, low signal-to-noise ratios, and thin liquidity in many
European venues.  This work package develops robust AI methods tailored to these
conditions.

\subsection*{Key Research Questions}
\begin{itemize}[nosep]
  \item How can transformer architectures be adapted for multi-asset financial
        time series with irregular timestamps?
  \item What reinforcement learning strategies are effective for portfolio
        optimization in low-liquidity European markets?
  \item How should model performance be evaluated when standard backtesting
        assumptions break down?
  \item Can transfer learning reduce data requirements for emerging-market
        financial instruments?
  \item What ensemble methods best balance predictive accuracy with robustness?
\end{itemize}

\subsection*{Methodological Approach}

Researchers will employ deep learning (LSTMs, transformers, attention
mechanisms), reinforcement learning (policy gradient, actor-critic), and
advanced time-series econometrics, validated on proprietary and public European
market data.

\subsection*{Expected Outcomes}
\begin{itemize}[nosep]
  \item Novel architectures for financial time-series prediction
  \item Open-source backtesting framework with realistic market friction modeling
  \item Empirical studies on AI-driven trading in European equity and bond
        markets
\end{itemize}

\subsection*{Related Training}
MST-02 (Introduction to AI for Financial Applications),
AST-05 (Reinforcement Learning in Digital Finance),
AST-06 (Machine Learning in Industry), AST-07 (Deep Learning for Finance),
AST-08 (Data-Centric AI).

\subsection*{Contributing Partners}
\begin{center}
\partnerlogo{WU}\hspace{3mm}%
\partnerlogo{ASE}\hspace{3mm}%
\partnerlogo{CAR}
\end{center}

\newpage

% ====================================================================
% WP 3
% ====================================================================
\wpheader{3}{Towards Explainable and Fair AI-Generated Decisions}{BFH Bern}

\subsection*{Research Motivation}

Researchers participating in WP3 will investigate how AI-based decision systems
in finance can be made transparent, interpretable, and fair.  As algorithmic
lending, insurance pricing, and credit scoring become widespread, regulators and
society demand that these systems explain their outputs and avoid
discriminatory bias.  The EU AI Act and forthcoming technical standards make this
work especially timely.

\subsection*{Key Research Questions}
\begin{itemize}[nosep]
  \item How can post-hoc explainability methods (SHAP, LIME, counterfactual
        explanations) be reliably applied to complex financial models?
  \item What fairness metrics are most appropriate for credit scoring and
        insurance pricing under European anti-discrimination law?
  \item How do explainability and accuracy trade off in practice, and can
        inherently interpretable models match black-box performance?
  \item What audit procedures will satisfy the EU AI Act's requirements for
        high-risk financial AI systems?
\end{itemize}

\subsection*{Methodological Approach}

Researchers will develop and benchmark explainability techniques on real-world
lending and insurance datasets, design fairness-aware learning algorithms, and
prototype compliance audit tools aligned with the EU AI Act.

\subsection*{Expected Outcomes}
\begin{itemize}[nosep]
  \item Comparative framework for XAI methods in financial decision-making
  \item Fairness-aware credit scoring models with auditable bias metrics
  \item Prototype EU AI Act compliance toolkit for financial institutions
  \item Best-practice guidelines for deploying explainable AI in regulated
        settings
\end{itemize}

\subsection*{Related Training}
MST-03 (The Need for eXplainable AI), MST-06 (Ethics Applicable to Digital Aspects),
AST-09 (Cybersecurity in Digital Finance), AST-10 (AI Design in Digital Finance),
AST-11 (Barriers in Digital Finance Adoption), AST-12 (Explainable AI in Finance).

\subsection*{Contributing Partners}
\begin{center}
\partnerlogo{BFH}\hspace{3mm}%
\partnerlogo{UTW}\hspace{3mm}%
\partnerlogo{UNA}
\end{center}

\newpage

% ====================================================================
% WP 4
% ====================================================================
\wpheader{4}{Driving Digital Innovation with Blockchain Applications}{Bucharest University of Economic Studies (ASE)}

\subsection*{Research Motivation}

Researchers participating in WP4 will explore how distributed ledger
technologies are reshaping financial infrastructure.  From decentralized
finance (DeFi) protocols and asset tokenization to central bank digital
currencies (CBDCs), blockchain introduces both transformative opportunities and
novel risks.  This work package examines the technical, economic, and regulatory
dimensions of blockchain-based financial innovation.

\subsection*{Key Research Questions}
\begin{itemize}[nosep]
  \item What smart-contract design patterns minimize security vulnerabilities in
        DeFi lending and trading protocols?
  \item How can real-world assets be tokenized efficiently and in compliance with
        European securities regulation?
  \item What CBDC architectures best balance privacy, programmability, and
        monetary policy transmission?
  \item How should regulatory sandboxes be designed to foster blockchain
        innovation while protecting consumers?
\end{itemize}

\subsection*{Methodological Approach}

Researchers will combine formal verification of smart contracts, agent-based
simulation of DeFi ecosystems, empirical analysis of on-chain data, and
comparative regulatory analysis of European sandbox frameworks.

\subsection*{Expected Outcomes}
\begin{itemize}[nosep]
  \item Secure smart-contract templates for tokenization and DeFi applications
  \item Agent-based models of DeFi market dynamics and systemic risk
  \item Policy briefs on CBDC design for European central banks
  \item Regulatory analysis of blockchain sandboxes across EU member states
\end{itemize}

\subsection*{Related Training}
MST-04 (Introduction to Blockchain Applications in Finance),
AST-13 (Digital Finance Regulation), AST-14 (History and Prospects of Digital Finance),
AST-15 (Blockchains in Digital Finance).

\subsection*{Contributing Partners}
\begin{center}
\partnerlogo{ASE}\hspace{3mm}%
\partnerlogo{UTW}\hspace{3mm}%
\partnerlogo{WU}\hspace{3mm}%
\partnerlogo{POZ}
\end{center}

\newpage

% ====================================================================
% WP 5
% ====================================================================
\wpheader{5}{Sustainability of Digital Finance}{University of Naples (UNA)}

\subsection*{Research Motivation}

Researchers participating in WP5 will address the dual challenge of making
finance more sustainable and making digital technologies themselves
environmentally responsible.  As the EU Taxonomy, SFDR, and CSRD reshape
disclosure requirements, financial institutions need robust tools to measure,
report, and act on climate risk and ESG performance.  At the same time, the
energy footprint of AI and blockchain systems demands attention.

\subsection*{Key Research Questions}
\begin{itemize}[nosep]
  \item How can NLP and large language models reliably extract ESG signals from
        unstructured corporate disclosures and news?
  \item What quantitative frameworks best integrate climate-risk scenarios into
        portfolio optimization and stress testing?
  \item How should green bond markets be evaluated for additionality and
        greenwashing risk?
  \item What metrics and practices can reduce the environmental impact of
        AI training and blockchain consensus mechanisms (``Green AI'')?
\end{itemize}

\subsection*{Methodological Approach}

Researchers will apply NLP-based ESG scoring, climate-scenario analysis
(NGFS scenarios), portfolio optimization under sustainability constraints, and
lifecycle assessment of computational workloads, drawing on data from consortium
banking partners and public ESG databases.

\subsection*{Expected Outcomes}
\begin{itemize}[nosep]
  \item NLP pipeline for automated ESG report analysis
  \item Climate-risk stress-testing framework for European bank portfolios
  \item Empirical assessment of green bond market integrity
  \item Green AI guidelines for energy-efficient model training in finance
\end{itemize}

\subsection*{Related Training}
MST-05 (Sustainable Finance), AST-16 (Digital EIT Summer School),
AST-17 (Green Digital Finance),
AST-18 (Multi-Criteria Decision Making in Sustainable Finance).

\subsection*{Contributing Partners}
\begin{center}
\partnerlogo{UNA}\hspace{3mm}%
\partnerlogo{UTW}\hspace{3mm}%
\partnerlogo{KUT}
\end{center}
